\chapter{Real-Time Integration Contract \& Operations}
\label{sec:deployment}

This chapter specifies the real-time integration surface and reports the
live run.  The stage is a \texttt{cudaPilotProxyDetector} stage in
\texttt{kotekan} \cite{kotekan}, the real-time framework CHIME and CHORD share,
and it is implemented for deployment on the CHORD Pathfinder for the reasons stated in
Section~\ref{sec:system:platforms}: CHIME's production pipeline is not a place
to try an unproven kernel, the Pathfinder runs the same framework with small
pipeline differences, and its commissioning state means an experimental stage
can run there without contending for science time or risking a production data
product.  The Pathfinder is therefore not a substitute for CHIME but the
telescope on which the live contract of this chapter --- framing, failure
semantics, exact replay, headroom, monitoring, and hours-long shadow operation
--- can actually be exercised.  Keeping the platform explicit avoids using
either an archive result or a synthetic framework harness as evidence for a
live system property, and it keeps the two calibrations apart: the archive
calibration of Chapter~\ref{sec:calibration} is CHIME's, and the Pathfinder
run establishes live behaviour on the Pathfinder's own receiver and
transmitter environment.  Table~\ref{tab:deployment:status} states where the stage stands, so that
``implemented'' is never read as ``deployed.''

\begin{table}[H]
\centering
\footnotesize
\caption[Status of the live stage]{Status of the live stage at the time of writing. Each row is a distinct milestone; a later row never implies an earlier one.}
\label{tab:deployment:status}
\begin{tabularx}{\linewidth}{>{\raggedright\arraybackslash}p{0.30\linewidth}>{\raggedright\arraybackslash}p{0.10\linewidth}>{\raggedright\arraybackslash}X}
\toprule
Milestone & Status & Evidence \\
\midrule
GPU kernel, CHIME profile & complete & Section~\ref{sec:impl:verify}: golden vectors, oracle, device equality, timing \\
GPU kernel, Pathfinder profile & complete & regenerated weights, bounds, and golden vectors (Section~\ref{sec:dtv:scope}) \\
\texttt{kotekan} host stage & complete & bundle binding, buffer contract, GPU-oracle and file-pipeline replay tests \\
Merged into the shared \texttt{kotekan} tree & \rerun{pending} & WVURAIL fork; upstream merge not required for the Pathfinder run \\
Shadow run on the Pathfinder & \pending & Section~\ref{sec:deployment:results} \\
Applied-mask trial & \pending & requires a passed shadow run and, on CHIME, the calibration bundles of Chapter~\ref{sec:calibration} \\
Science-transfer test & \pending & contaminated visibilities before and after the sidereal filter (Section~\ref{sec:tolerance:chain}) \\
\bottomrule
\end{tabularx}
\end{table}

\section{Stage Boundary and Frame Alignment}
\label{sec:deployment:stage}

The production stage sits after channelization and before visibility
accumulation.  It consumes the same packed samples that contribute to a
visibility integration and emits one validity/mask decision for the
corresponding monitored DTV allocation.  The acceptance condition is integer
alignment: every detector frame must cover exactly one visibility integration
or an explicitly declared integer number of them.  A partial overlap is
rejected at configuration time because it has no unambiguous mask semantics.

The Pathfinder stage realizes that placement.  Its input is the packed
voltage ring with logical shape $[T,F,P,D]$ (time, frequency, polarization,
dish); metadata bind each local coarse channel to the runtime bundle's
\texttt{chord\_channel\_id}.  At $f_s = 195.3125$~kHz it uses the Pathfinder
profile of Section~\ref{sec:dtv:scope}, $K=64$ and $N=8192$, or 128 windows
spanning $41.94304$~ms, exactly one configured Pathfinder visibility
integration.  It writes one \texttt{int8} mask and three \texttt{uint64} coarse
marginals per local frequency into a downstream buffer, and the mask is
consumed by the visibility stage of the same frame or ignored, by
configuration, in shadow mode.  Figure~\ref{fig:deployment:pipeline} places the
stage in the pipeline.  This is concrete alignment evidence for the Pathfinder
geometry; the CHIME geometry follows the same placement with its own profile.

\begin{figure}[t]
\centering
\resizebox{\linewidth}{!}{% Editable source for Figure \ref{fig:deployment:pipeline}: the stage in the kotekan pipeline.
\begin{tikzpicture}[font=\small, x=1cm, y=1cm]
\node[blk, text width=22mm, minimum height=10mm] (fe) at (0,0) {F-engine /\\network receive};
\pic at (3.0,0) {ring};
\node[lab, text width=26mm] at (3.0,-0.95) {packed voltage ring $[T,F,P,D]$, \texttt{int4}};
\node[blk, draw=ppMeasured!90!black, fill=ppMeasured!12, line width=1pt, text width=48mm, minimum height=12mm] (det) at (7.4,0) {\textbf{pilot-proxy stage} (\texttt{cudaPilotProxyDetector})\\one GPU block per input; runtime bundle: channel ids, weights, $(f_a,\mathcal D,\mathcal B,\rho,\eta_{q16})$};
\pic at (6.0,-1.9) {threads};
\node[blk, text width=38mm, minimum height=10mm] (out) at (12.2,0) {output buffer, per frame and channel: mask $+$ validity (\texttt{int8}), $3\times$\texttt{uint64} marginals};
\node[blk, text width=22mm, minimum height=10mm] (xe) at (16.4,0) {visibility accumulation};
\node[blk, text width=22mm, minimum height=9mm] (down) at (16.4,-2.6) {downstream consumers};
\node[blkx, text width=44mm, minimum height=11mm] (mon) at (7.4,-3.7) {monitoring and shadow record: tripwire and validity reasons, runtime quantiles, drift, sampled frames for exact replay};
\draw[bus] (fe) -- (2.55,0);
\draw[bus] (3.45,0) -- (det) node[width, pos=0.72, above]{$M\times N$ int4};
\draw[sig] (det) -- (out); \draw[sig] (out) -- (xe);
\draw[bus] (3.0,0.45) -- ++(0,1.4) -| (xe.north) node[lab, pos=0.25, above]{same samples, same frame};
\draw[sigc] ([yshift=-3mm]out.east) -- ++(0.55,0) |- (down.west);
\draw[ctl] (det.south) -- (mon.north);
\node[lab, anchor=west, text width=26mm] at (10.0,-3.7) {$\leftarrow$ replayed through the offline reference: exact agreement required};
\node[lab, text width=36mm] at (12.2,-1.65) {shadow mode: the visibility stage ignores the mask; applied mode: consumes it for the same frame};
\end{tikzpicture}
}
\caption[The stage in the \texttt{kotekan} pipeline]{The stage in the \texttt{kotekan} pipeline. It reads the packed voltage ring that the visibility stage also reads, evaluates every monitored coarse channel once per frame from a runtime bundle, and writes one mask-and-validity byte and three exact coarse marginals per channel per frame. In shadow mode the visibility stage ignores the mask; in applied mode it consumes it for the same frame. The monitoring surface records the tripwire and validity state, runtime quantiles, drift indicators, and sampled frames for exact replay against the offline reference.}
\label{fig:deployment:pipeline}
\end{figure}

\begin{figure}[t]
\centering
\resizebox{\linewidth}{!}{% Editable source for Figure \ref{fig:deployment:lifecycle}: the stage lifecycle as a decision flow.
\begin{tikzpicture}[font=\footnotesize, node distance=6mm and 8mm,
  st/.style={draw, rounded corners=2pt, align=center, inner sep=4pt, minimum height=15mm, text width=#1, line width=0.75pt},
  ex/.style={st=#1, fill=ppMeasured!8, draw=ppMeasured!90!black},
  fail/.style={st=#1, fill=ppFailure!8, draw=ppFailure!90!black},
  ok/.style={st=#1, fill=ppConditional!10, draw=ppConditional!90!black},
  dec/.style={diamond, draw=ppModel!90!black, fill=ppModel!10, align=center, inner sep=1pt, aspect=2.4, font=\footnotesize, line width=0.75pt},
  arr/.style={-{Stealth[length=1.8mm]}, thick, ppMeasured!90!black},
  arrf/.style={-{Stealth[length=1.8mm]}, thick, ppFailure!90!black}]
\node[ex=32mm] (in) {\textbf{input frame}\\frame identity, geometry, era, receiver contract};
\node[dec, right=of in] (c1) {contracts\\valid?};
\node[ex=34mm, right=of c1] (init) {\textbf{initialize}\\bind the bundle by era; zero accumulators and the completion word, in-stream};
\node[ex=34mm, right=of init] (launch) {\textbf{launch}\\one block per input: row sums, marginals, \texttt{fxfft256}, atomic sums};
\node[ex=34mm, below=9mm of launch] (fin) {\textbf{finalize} (last block)\\fence; re-read the sums; marginal tripwire; exact rank and threshold};
\node[dec, left=of fin] (c2) {CUDA ok and\\tripwire exact?};
\node[ok=34mm, left=of c2] (pub) {\textbf{publish valid}\\mask bit, valid$=$true, marginals, diagnostics; completion event};
\node[fail=64mm, below=9mm of c2, xshift=-8mm] (bad) {\textbf{publish invalid}\\valid$=$false with a reason code, completion event, operations event --- never substitute keep or mask; an invalid frame enters neither retained nor excised accounting};
\draw[arr] (in) -- (c1);
\draw[arr] (c1) -- node[above, font=\scriptsize]{yes} (init);
\draw[arr] (init) -- (launch);
\draw[arr] (launch) -- (fin);
\draw[arr] (fin) -- (c2);
\draw[arr] (c2) -- node[above, font=\scriptsize]{yes} (pub);
\draw[arrf] (c1.south) -- ++(0,-4mm) node[right, font=\scriptsize]{no} -| ([xshift=-5mm]pub.west) |- (bad.west);
\draw[arrf] (c2.south) -- node[right, font=\scriptsize]{no} (bad.north -| c2.south);
\end{tikzpicture}
}
\caption[Real-time stage lifecycle and failure semantics]{The required production
stage lifecycle.  Initialization is part of correctness: geometry and epoch
validity are checked before launch, outputs are zeroed in the same CUDA stream,
and a CUDA, completion-counter, or exact-tripwire failure produces an invalid
frame and operations event.  It never silently substitutes ``keep'' or ``mask''
for an execution whose decision contract was not satisfied.  The Pathfinder stage
exposes this lifecycle, including the separate reason-coded validity surface
shown here, and the shadow run is where it is exercised under real operations.}
\label{fig:deployment:lifecycle}
\end{figure}

\begin{table}[t]
\centering
\caption[Real-time interface contract]{Real-time interface contract.  Concrete C++ type names may change with
the host framework; the semantic fields and failure outcomes may not.}
\label{tab:deployment:interface}
\small
\begin{tabularx}{\linewidth}{>{\raggedright\arraybackslash}p{0.18\linewidth}
                             >{\raggedright\arraybackslash}p{0.30\linewidth}X}
\toprule
Surface & Required contract & Failure outcome \\
\midrule
Input samples & Contiguous packed complex \texttt{int4}; declared $M,N$, channel, frame identity, and time; immutable until stream completion & Frame invalid; no launch on size or geometry mismatch \\
Input health & Reason-coded constant/fill, code-histogram, variance, and saturation checks independent of the detection statistic & Frame invalid and quarantined for input-path diagnosis \\
Calibration & Channel and epoch validity; weight/norm/twiddle hashes; anchor, masks, rank, Q16 multiplier & Bundle rejected before data taking \\
CUDA stream & Input readiness event, stream-ordered output initialization, launch, completion event, and downstream wait & Frame invalid and stage health counter incremented \\
Fine/coarse outputs & Declared spans for power terms, diagnostics, completion/mask word, and optional debug taps & No truncation; size mismatch is configuration failure \\
Mask semantics & One aligned decision plus a separate validity bit; invalid is not encoded as clean & Downstream policy decides whether invalid data are retained, quarantined, or stop the run \\
Provenance & Source, binary, toolkit, schema, and configuration hashes; epoch identifier attached to monitoring record & Run cannot enter science mode without complete manifest \\
\bottomrule
\end{tabularx}
\end{table}

Figure~\ref{fig:deployment:timing} draws the same contract on a time axis, which is where the ``same frame'' question is actually settled: the detector can only run once the ring holds all of frame $n$, its kernel finishes in about a millisecond of the $41.94$~ms period, and in applied mode the visibility stage's integration of frame $n$ waits on the completion event that publishes mask$(n)$ and valid$(n)$; in shadow mode it does not wait, and the mask is recorded beside the integration it would have gated. Either way the deadline is the same: everything for frame $n$ is done before frame $n{+}1$ is ready.

\begin{figure}[t]
\centering
\resizebox{\linewidth}{!}{% Editable source for Figure \ref{fig:deployment:timing}: one frame on a time axis (tikz-timing).
\begin{tikztimingtable}[timing/slope=0.18, timing/xunit=0.78cm, timing/yunit=0.6cm, timing/rowdist=5ex, timing/coldist=2ex,
  timing/font=\scriptsize, timing/name/.style={font=\scriptsize, align=right}]
frame clock (rising edge $=$ readiness) & 6L 3H 3L 3H 3L \\
samples in the ring & 6D{frame $n-1$} 6D{frame $n$} 6D{frame $n+1$} \\
detector kernel (busy) & 6L 0.15L 0.6H 5.25L 0.15L 0.6H 5.25L \\
completion event (mask $+$ validity) & 6.75L 0.2H 5.8L 0.2H 5.05L \\
visibility, applied mode & 6.95L 5.05D{frame $n$, masked} 0.95L 5.05D{frame $n+1$, masked} \\
visibility, shadow mode & 6D{frame $n-1$} 6D{frame $n$, unmasked; mask($n$) logged} 6D{frame $n+1$} \\
\extracode
\begin{pgfonlayer}{background}
\vertlines[ppPending!60, dashed]{6, 12, 18}
\end{pgfonlayer}
% timing annotations, in timing units
\draw[<->, ppMeasured!90!black, thick] (6,1.7) -- (12,1.7) node[midway, above, font=\scriptsize, text=ppMeasured!90!black]{$t_{\rm frame} = N/f_s = 41.94$~ms};
\draw[<->, ppFailure!90!black, thick] (6.15,-7.1) -- (6.95,-7.1); \node[font=\scriptsize, text=ppFailure!90!black, anchor=east, align=right] at (6.0,-7.1) {$t_{\rm kernel} \approx 1$~ms (drawn wider)};
\draw[<->, ppConditional!90!black, thick] (6.95,-8.1) -- (12,-8.1) node[midway, below, font=\scriptsize, text=ppConditional!90!black]{headroom: everything for frame $n$ before frame $n{+}1$ is ready (the deadline)};
\node[font=\scriptsize, text=ppMeasured!90!black, anchor=north, text width=120mm, align=center] at (9,-10.3) {the frame clock's rising edge is the readiness event; the completion event is the only other signal in the contract: applied-mode integration waits on it, shadow-mode integration does not};
\end{tikztimingtable}
}
\caption[Frame timing: readiness, launch, completion, deadline]{What happens to one frame, and when. Samples stream continuously; a readiness event marks frame $n$ complete in the ring; the detector launches on it and finishes in about a millisecond; a completion event publishes the mask and validity for frame $n$; in applied mode the visibility integration of frame $n$ waits on that event, in shadow mode it does not; and the downstream deadline is that all of frame $n$ is finished before frame $n{+}1$ is ready.}
\label{fig:deployment:timing}
\end{figure}

Two geometry statements must not be conflated.  The Pathfinder stage uses
$N=8192$, $K=64$, and $f_s=195.3125$~kHz and declares one-to-one alignment
with its $8192$-sample visibility integration.  The benchmarked CHIME detector
uses $N=16384$, $K=128$, and $f_s=390.625$~kHz for the same 41.94-ms detector
cadence --- the upgraded-engine frame, whereas the published CHIME correlator
integrates over a different length.  A CHIME live stage therefore passes only
after that engine's actual frame contract is pinned and one-to-one or
integer-to-one alignment is demonstrated with injected frame identities.
Regenerating for another power-of-two geometry creates new weights, transform
bounds, golden vectors, and performance acceptance; it is not an informal
runtime switch, and it is exactly what the Pathfinder port did.

\section{Initialization, Synchronization, and Failure Semantics}
\label{sec:deployment:failure}

The host wrapper owns all state that the kernel assumes.  Before each launch it
zeros the fine and coarse accumulation buffers and the completion/mask word in
the same CUDA stream, binds a calibration bundle whose epoch includes the frame
time, and records the input-ready dependency.  The downstream consumer waits on
the detector completion event before reading either the mask or validity bit.
The raw kernel entry point remains internal so a caller cannot omit the
completion-word initialization described in Chapter~\ref{sec:impl}.

A live system must distinguish a scientific non-detection from an execution
failure.  The following conditions set \texttt{valid=false}: geometry mismatch,
constant/fill or saturation health failure, invalid or expired calibration,
zero-reference/bulk degeneracy beyond the
specified detector rules, CUDA launch or asynchronous error, impossible
completion count, exact-marginal tripwire failure, output checksum failure, or
missing provenance.  The stage records the reason and the affected frame range.
It does not manufacture a clean decision.  The collaboration's downstream data
policy may conservatively quarantine invalid integrations, but that policy is
outside the detector bit and remains visible.

\section{Operations Runbook}
\label{sec:deployment:runbook}

A concise run sequence is sufficient because the calibration and binary are
content-addressed artifacts rather than editable files:

\begin{enumerate}[leftmargin=2em]
  \item \textbf{Pre-run validation.} Verify receiver geometry, frame alignment,
  GPU capability, free device memory, bundle epoch, all hashes, and the signed
  pilot-injection convention.  Reject mixed or expired artifacts.
  \item \textbf{Acceptance burst.} Process a deterministic quiet vector, a
  signed-offset injected pilot, a constant detector-input packed-\texttt{0x88}
  input-health failure (native excess-8 \texttt{0x00} before repacking), a
  zero-denominator case, and a nonzero-prefilled mask-word case.
  Require equality with the independent reference and the staged kernel path.
  \item \textbf{Science mode.} Enable the stage only after the acceptance burst;
  record per-frame validity, mask, coarse/fine tripwire status, kernel duration,
  and calibration identifier.
  \item \textbf{Respond to a tripwire.} Stop promoting affected frames to science
  products, preserve the input/product identifiers, capture a debug-tap
  reproduction, and compare against the reference before resuming.
  \item \textbf{Change control.} Treat any source, compiler, toolkit, geometry,
  twiddle, weight, schema, or calibration change as a new manifest.  Code or
  geometry changes repeat the full verification suite; an epoch calibration
  change repeats the calibration and acceptance gates.
\end{enumerate}

\section{Monitoring: Tripwire, Sentinel, and Drift}
\label{sec:deployment:monitoring}

The exact-marginal tripwire compares independently accumulated coarse and fine
marginals as integers.  It detects implementation and memory-path disagreement;
it does not detect an astrophysical line at the wrong frequency.  A separate
schedule-agnostic sentinel compares persistent coarse excess with the
fine-designated decision.  Sustained coarse excess on frames the designated
window does not reject is evidence for a moved anchor, an additional transmitter,
or power outside the current census mask and triggers re-localization at the
next epoch boundary.

This sentinel is required by a concrete legacy failure, not only by defensive
design.  In the earlier 23-channel survey, channel 33 appeared healthy inside
the statistic --- a near-$0.5$ mask fraction and no measured target-cell leak ---
while its extracted carriers lay in the skipped guard near $-3.6$~kHz with
approximately $21$~dB integrated SNR.  The detector response there was about
$-17$~dB.  A strong off-nominal pilot could therefore be less visible than a
weaker nominal one: ``clean'' and ``blind'' were indistinguishable from the
target statistic alone.  The same survey separated five likely instrument
lines by rational receiver-clock coincidence, single-channel localization, and
mask invariance, while retaining an unidentified sixth line rather than
excising it by assumption.  Appendix~\ref{app:legacy} records those diagnostics
and their unresolved attribution.  Operationally, the lesson is that a target
statistic must be monitored by a broader spectral sentinel and a periodically
refreshed line census.

Monitoring therefore has three layers: exact arithmetic health, statistical
null/false-alarm health on non-designated bins, and environment drift.  Dashboard
rates are aggregated by channel, epoch, GPU, and calibration hash so an apparent
broadcast change is not confused with a software rollout.  The live integration
is complete only when these records, the invalid-frame handling, and a failure
injection exercise have been demonstrated on the target telescope's actual
\texttt{kotekan} buffer and downstream paths.  The file-pipeline parity test
is a prerequisite for that exercise, not a substitute for it.

\section{Performance Acceptance}
\label{sec:deployment:performance}

On an A100 benchmark GPU, the staged samples-to-powers path is reported at
\rerun{2.14~ms}, the fused path at \rerun{0.77~ms}, and the fused path with
decision epilogue at \rerun{0.90~ms} for the $N=16384$ geometry
(Section~\ref{sec:impl:verify}).  These measurements demonstrate substantial
compute margin relative to a 41.94-ms detector frame.  They do not by themselves
include framework scheduling, input readiness, inter-stage contention, or
visibility-buffer synchronization, and they are not a full-load measurement of
the $N=8192$ Pathfinder stage.  The live acceptance test therefore reports
end-to-end latency percentiles and deadline misses under the full telescope load,
with the kernel timing retained as a component rather than presented as the
stage latency.
\stub{Pathfinder shadow run}{Full-load numbers: per-frame stage latency
percentiles (50/90/99/max) on the Pathfinder GPU, deadline misses, memory
headroom, and the kernel-only time on that GPU for the $N=8192$ profile.}

\section{The Pathfinder Shadow Run}
\label{sec:deployment:results}

The closing live measurement is a predeclared shadow run of at least several
hours on the CHORD Pathfinder, with one frozen runtime bundle per monitored
allocation and the mask recorded but not applied.  It is the online level of
both evidence ladders at once (Sections~\ref{sec:estimator:runtime}
and~\ref{sec:detection:runtime}), and it reports five things.

\begin{enumerate}[leftmargin=1.6em]
\item \emph{Framing and stability.} Frames processed, dropped, and late per
channel; validity reasons; tripwire state; anchor and bulk drift by time block.
\item \emph{Exact replay.} Sampled live frames replayed through the
independent offline reference, compared as integers: coarse marginals, fine
powers, and the mask bit.  The expected mismatch count is zero.
\item \emph{Headroom.} The latency and memory figures of
Section~\ref{sec:deployment:performance} under real load.
\item \emph{What the mask removes.} For each monitored channel, the
frame-averaged power spectrum of all frames and of the frames the mask keeps,
accumulated over the run --- the live counterpart of the archive plates of
Appendix~\ref{app:archive-diagnostics} --- together with the mask fraction as
a time series.  On the Pathfinder these spectra describe the Pathfinder's
transmitter environment; they are a live demonstration of the detector, not a
CHIME channel verdict.
\item \emph{Response to state changes.} Where a transmitter sign-on, sign-off,
or a deliberate bundle swap occurs during the run, the recorded response
against the declared one.
\end{enumerate}

The acceptance criteria are numerical and are fixed before the run
(Table~\ref{tab:deployment:acceptance}); the blue values are the proposed
settings and are to be confirmed or replaced before the run is scheduled, not
after it is read.

\begin{table}[H]
\centering
\footnotesize
\caption[Shadow-run acceptance criteria]{Acceptance criteria for the Pathfinder shadow run, fixed before the run. Failing any row fails the run; a failed run is reported with the failing row, not silently repeated.}
\label{tab:deployment:acceptance}
\begin{tabularx}{\linewidth}{>{\raggedright\arraybackslash}p{0.30\linewidth}>{\raggedright\arraybackslash}p{0.22\linewidth}>{\raggedright\arraybackslash}X}
\toprule
Criterion & Requirement & Basis \\
\midrule
Duration and frames & $\ge$ \rerun{24}~h continuous ($\ge$ \rerun{$2.06\times10^6$} frames per monitored channel at $41.94$~ms); an \rerun{8}-h segment is the minimum engineering-stability run and does not support a diurnal claim & a diurnal propagation cycle needs a full day \\
Exact replay & \rerun{$10^4$} sampled frames per channel, uniformly over the run, zero mismatches in coarse marginals, fine powers, and mask bit & a single mismatch is a contract violation \\
Frame loss & dropped $+$ late $\le$ \rerun{$10^{-5}$} of frames, with no run of more than \rerun{2} consecutive losses & losses are logged with cause \\
Latency and memory & stage $p_{99.9}$ $\le$ \rerun{$0.5$} of the frame period and the maximum never exceeds one period; GPU memory headroom $\ge$ \rerun{20\%} under concurrent-stage contention & headroom against the $41.94$~ms deadline \\
Invalid frames & system failures (CUDA, tripwire, buffer) and scientifically invalid inputs (all-zero, railed, degenerate reference) counted separately by reason; any system-failure reason exceeding \rerun{$10^{-4}$} of frames, or any input reason exceeding \rerun{$10^{-3}$}, halts the run for diagnosis & validity is a product, not a filter \\
Tripwire & zero marginal-identity failures & exact contract \\
Drift & anchor and bulk drift within the era's calibration interval for every time block & Section~\ref{sec:deployment:monitoring} \\
Injected failures & buffer starvation, bundle swap, deliberate all-zero and railed frames, and a forced kernel fault each produce the declared response (invalid frame, event, no silent keep/mask) & Section~\ref{sec:deployment:failure} \\
\bottomrule
\end{tabularx}
\end{table}

\stubfig{Pathfinder shadow run}{One representative channel: the run's
all-frame spectrum and the \emph{counterfactual} kept-frame spectrum (the mask
is recorded, not applied, so ``after'' is what the mask would have kept), on
the archive-plate axes; beside it the run-health timeline (mask fraction,
validity, tripwire, latency). The remaining channels' panels go to an
appendix.}{All-frame and counterfactual kept-frame spectra with the run-health
timeline, one representative channel of the Pathfinder shadow
run.\label{fig:deployment:spectra}}

\stubtab{Pathfinder shadow run}{Two tables. Run-global: duration, frames
processed, dropped and late, latency percentiles and maximum, memory headroom,
exact-replay frames compared and mismatches, system-failure counts by reason.
Per channel: validity counts by input reason, tripwire firings, mask fraction
and its range over time blocks, anchor and bulk drift, state changes observed
and the stage's response, disposition.}{Shadow-run summary: run-global health
and per-channel behaviour on the CHORD Pathfinder.\label{tab:deployment:results}}

Only a successful shadow run can authorize a later applied-mask trial, and an
applied-mask trial on CHIME additionally requires the CHIME calibration
bundles of Chapter~\ref{sec:calibration} and the blocked holdout of
Section~\ref{sec:survey:holdout}; nothing in the Pathfinder run substitutes for
either.
